\documentclass[final,5p,times,authoryear]{elsarticle}

\usepackage[utf8]{inputenc}
\usepackage[T1]{fontenc}
\usepackage[spanish,es-nodecimaldot]{babel}
\usepackage{graphicx}
\usepackage{booktabs}
\usepackage{tabularx}
\usepackage{array}
\usepackage{multirow}
\usepackage{adjustbox}
\usepackage{longtable}
\usepackage{url}
\usepackage{xcolor}
\usepackage{hyperref}
\usepackage{placeins}

\addto\captionsspanish{\renewcommand{\tablename}{Tabla}}

\makeatletter
\pdfstringdefDisableCommands{%
  \def\@fnmark{}%
  \def\fnref#1{}%
  \def\corref#1{}%
  \def\tnoteref#1{}%
  \def\comma{, }%
}
\makeatother

\hypersetup{
  colorlinks=true,
  linkcolor=black,
  citecolor=black,
  urlcolor=blue,
  pdftitle={Análisis computacional de salud mental en Reddit: estado del arte en PLN y modelos de lenguaje},
  pdfauthor={Jaramillo P. Ulises; Ponce M. Fernanda; Vivas M. Julio C.}
}

\graphicspath{{figures/}}

\begin{document}

\begin{frontmatter}

\title{Análisis computacional de salud mental en Reddit: estado del arte en PLN y modelos de lenguaje}

\author[tec]{Jaramillo P. Ulises\fnref{ulises}}
\author[tec]{Ponce M. Fernanda\fnref{fernanda}}
\author[tec]{Vivas M. Julio C.\fnref{julio}}

\fntext[ulises]{Matrícula: A01798380}
\fntext[fernanda]{Matrícula: A01799293}
\fntext[julio]{Matrícula: A01749879}

\address[tec]{TC3002B.501 Desarrollo de aplicaciones avanzadas de ciencias computacionales, Equipo 6, Tecnológico de Monterrey Campus Estado de México, Escuela de Ingeniería y Ciencias}

\begin{abstract}
Este manuscrito corresponde al entregable de Fase 1 del reto TC3002B.501 y presenta una revisión del estado del arte sobre la detección computacional de señales de salud mental en texto de redes sociales, con énfasis en Reddit. En lugar de proponer un modelo experimental, este trabajo se concentra en analizar críticamente literatura científica reciente para identificar tendencias, aportaciones, limitaciones y oportunidades de investigación relevantes para el desarrollo posterior de una herramienta computacional. La revisión integra trece artículos sobre procesamiento de lenguaje natural (PLN), aprendizaje automático, aprendizaje profundo, modelos basados en transformers y modelos de lenguaje de gran escala aplicados a depresión, ideación suicida y enfoques clínicamente fundamentados en DSM-5. El análisis se organiza a partir de la dimensión seleccionada para el proyecto, a saber, la efectividad predictiva y la interpretabilidad clínica. A partir de esta lectura comparativa, se distinguen avances importantes en desempeño cuantitativo, así como brechas persistentes en explicabilidad, anotación experta, generalización y utilidad clínica.
\end{abstract}

\begin{keyword}
salud mental \sep redes sociales \sep Reddit \sep procesamiento de lenguaje natural \sep aprendizaje automático \sep detección de depresión \sep ideación suicida \sep DSM-5 \sep modelos de lenguaje
\end{keyword}

\end{frontmatter}

\section{Introducción}

El análisis computacional de salud mental en texto generado por usuarios se ha consolidado como una línea de investigación estratégica para la detección temprana de señales de riesgo, especialmente en entornos donde la búsqueda de ayuda profesional se retrasa por estigma, costo o falta de acceso. En ese contexto, las plataformas de redes sociales ofrecen narrativas espontáneas, longitudinales y, en muchos casos, suficientemente detalladas para estudiar indicadores de depresión, desesperanza, aislamiento, autolesión e ideación suicida \citep{montejo2024survey, abdulsalam2024suicidal}.

Dentro de ese ecosistema, Reddit ocupa una posición particularmente relevante por su estructura basada en comunidades temáticas, la extensión relativamente mayor de sus publicaciones y la posibilidad de anonimato, factores que favorecen expresiones más explícitas sobre malestar psicológico. Trabajos tempranos y recientes muestran que el valor de Reddit no se limita a la clasificación binaria de riesgo, sino que también permite modelar transiciones hacia ideación suicida, niveles de severidad, categorías diagnósticas clínicamente significativas y explicaciones textuales de apoyo a la decisión \citep{dechoudhury2016shifts, gaur2019severity, bao2025redsm5}.

Este documento presenta una revisión de estado del arte, no un artículo experimental. Su objetivo es identificar y analizar literatura reciente y representativa sobre detección computacional de señales de salud mental en redes sociales, con énfasis en depresión, ideación suicida, marcos DSM-5, aprendizaje automático, aprendizaje profundo, modelos basados en transformers y modelos de lenguaje de gran escala. La revisión se organiza en función de la dimensión de desempeño seleccionada para el reto: \textit{efectividad predictiva e interpretabilidad clínica}. Esta elección permite evaluar no solo qué tan bien predicen los modelos, sino también si sus decisiones pueden vincularse con síntomas, criterios o evidencias comprensibles para profesionales de salud mental.

La hipótesis que articula la revisión es que el campo ha avanzado con rapidez en capacidad predictiva, pero no de manera uniforme en fundamentación clínica, transparencia y utilidad para intervención temprana. Por ello, el documento compara enfoques centrados en características lingüísticas, modelos profundos, recursos clínicamente anotados y sistemas basados en transformers y LLMs, destacando sus aportaciones, limitaciones y el tipo de evidencia que ofrecen para fases posteriores del proyecto.

\section{Contexto del problema y motivación}

La detección automática de señales de salud mental en redes sociales responde a una necesidad dual. Por una parte, existe una motivación de salud pública: la ideación suicida y la depresión requieren identificación oportuna para facilitar derivación, monitoreo o intervención. Por otra parte, existe una motivación computacional: convertir lenguaje informal, ruidoso y altamente contextual en representaciones útiles para estimar riesgo clínico de forma robusta y reproducible \citep{abdulsalam2024suicidal, montejo2024survey}.

El problema es especialmente complejo por tres razones. Primero, los textos de redes sociales no equivalen a expedientes clínicos: contienen ironía, jerga, referencias implícitas y variación discursiva significativa. Segundo, la clasificación útil para salud mental no siempre es binaria; en muchos casos importa distinguir cambios de estado, niveles de severidad, síntomas específicos o racionales textuales que expliquen la predicción \citep{dechoudhury2016shifts, gaur2019severity}. Tercero, los modelos con mejor desempeño cuantitativo suelen ser también los menos transparentes, lo que introduce un conflicto entre exactitud estadística y confianza clínica.

En Reddit, estas tensiones se vuelven particularmente visibles. Los subreddits de apoyo permiten estudiar desde discursos generales de malestar hasta comunidades explícitamente asociadas con suicidio, depresión, ansiedad o trastornos de personalidad. Ello ha permitido construir tareas distintas: detección de publicaciones suicidas, pronóstico de transición hacia \textit{r/SuicideWatch}, clasificación de severidad, mapeo a categorías DSM-5 y análisis explicativo con modelos de lenguaje \citep{gaur2018dsm5, gaur2019severity, bauer2024llm}.

Para este proyecto, el interés no se centra en cualquier aproximación de PLN, sino en aquellas que aportan evidencia para dos capacidades concretas. La primera es la \textbf{efectividad predictiva}, entendida como la capacidad de discriminar entre clases de riesgo o severidad con métricas sólidas como exactitud, precisión, recall, F1, AUC o variantes ordinales. La segunda es la \textbf{interpretabilidad clínica}, entendida como la posibilidad de relacionar la decisión del sistema con síntomas, categorías diagnósticas, racionales textuales o estructuras conceptuales útiles para especialistas. Esta segunda capacidad es crítica porque el reto académico no busca únicamente maximizar una métrica, sino justificar técnicamente por qué un enfoque podría ser pertinente para un dominio tan sensible.

\section{Método de selección de literatura}

La revisión se construyó a partir de un corpus curado de trece artículos científicos disponibles localmente en la carpeta \path{context/papers/}. El corpus incluye artículos de revista y conferencias publicadas entre 2016 y 2025, y fue delimitado para cubrir cuatro focos complementarios: detección general de trastornos mentales con PLN, detección de ideación suicida en redes sociales, enfoques de depresión fundamentados clínicamente en DSM-5 y desarrollos recientes basados en aprendizaje profundo, modelos basados en transformers y LLMs.

Los criterios de inclusión fueron los siguientes: (i) abordar explícitamente salud mental, depresión o ideación suicida en texto de redes sociales; (ii) utilizar o discutir metodologías de PLN, aprendizaje automático, aprendizaje profundo, modelos basados en transformers o LLMs; (iii) tener relevancia directa para Reddit o para tareas comparables de detección en redes sociales; y (iv) aportar evidencia para la dimensión de análisis seleccionada, es decir, efectividad predictiva y/o interpretabilidad clínica. Bajo estos criterios, se privilegiaron artículos que describen el origen de los datos, la representación textual, el modelo computacional, las métricas reportadas y las implicaciones prácticas de sus resultados.

La literatura seleccionada se analizó de forma comparativa y no como una lista de resúmenes independientes. Para ello, cada artículo fue caracterizado mediante los siguientes atributos: problema de salud mental abordado, fuente de datos, representación textual, modelo, métricas reportadas, contribución principal, limitaciones y aportación específica a la efectividad predictiva o a la interpretabilidad clínica. Este esquema permite alinear la revisión con la rúbrica de evaluación del curso, porque fuerza a explicitar no solo qué hace cada trabajo, sino también por qué su contribución resulta pertinente para el reto.

Finalmente, los trabajos fueron agrupados por enfoque de investigación. Esta organización facilita distinguir entre estudios panorámicos, modelos tempranos basados en rasgos lingüísticos, enfoques clínicamente anclados en DSM-5 y modelos contemporáneos que priorizan potencia predictiva mediante arquitecturas profundas o LLMs. La Tabla~\ref{tab:state-of-the-art} sintetiza el corpus completo y funciona como referencia transversal para las secciones analíticas posteriores.

\section{Estado del arte}

\subsection{Enfoques generales de PLN para detección de salud mental}

Los trabajos de revisión recientes muestran que la detección de salud mental mediante PLN ha evolucionado desde representaciones manuales basadas en palabras, léxicos y metadatos hacia arquitecturas profundas y modelos basados en transformers, pero con una persistencia notable de problemas estructurales. \citet{montejo2024survey} sintetizan tendencias, datasets, algoritmos y campañas de evaluación en varios trastornos, y subrayan carencias recurrentes: predominio del inglés, necesidad de anotación experta, bajo énfasis en explicabilidad y escasa validación en escenarios reales. Desde una perspectiva más específica, \citet{abdulsalam2024suicidal} revisan treinta y siete estudios sobre ideación suicida y confirman que el campo sigue concentrado en Twitter, Reddit y Weibo, con una mezcla de aprendizaje supervisado clásico, CNN, LSTM y modelos híbridos.

Estas revisiones son importantes porque desplazan la discusión desde el rendimiento aislado de cada artículo hacia la madurez metodológica del área. En términos de efectividad predictiva, ambos estudios muestran que el rendimiento depende no solo del modelo, sino también del esquema de anotación, el idioma, la plataforma y la granularidad de la tarea. En términos de interpretabilidad clínica, ambos coinciden en que la literatura tiende a operacionalizar el problema como clasificación textual y no como apoyo explicativo a la toma de decisiones, lo que deja una brecha entre desempeño computacional y utilidad clínica \citep{montejo2024survey, abdulsalam2024suicidal}.

Un antecedente clave para comprender esta tensión es el trabajo que modela el desplazamiento desde discurso general de salud mental hacia discurso suicida en Reddit, en lugar de limitarse a clasificar publicaciones aisladas \citep{dechoudhury2016shifts}. Su contribución es relevante porque introduce una lógica longitudinal y conductual: el riesgo no se expresa únicamente por palabras explícitas, sino por cambios en cohesión lingüística, foco en uno mismo, interacción social y marcadores de desesperanza. Desde la dimensión de interpretabilidad clínica, ese cambio de perspectiva aumenta el valor del análisis, ya que permite pensar en detección temprana y no solo en clasificación retrospectiva. Este trabajo aporta al reto porque muestra que la utilidad del sistema depende también de capturar trayectorias discursivas y no únicamente etiquetas finales.

\subsection{Detección de ideación suicida en redes sociales}

En la línea de clasificación de publicaciones suicidas, \citet{tadesse2020suicide} representan una etapa de transición entre aprendizaje automático convencional y aprendizaje profundo. Su combinación LSTM-CNN con \textit{word embeddings} mejora sobre modelos aislados y sobre clasificadores clásicos, mostrando que las arquitecturas híbridas capturan simultáneamente dependencias secuenciales y patrones locales. Desde la dimensión de efectividad predictiva, el trabajo es importante porque legitima el uso de redes profundas sobre Reddit; sin embargo, su interpretabilidad es reducida y los propios autores reconocen limitaciones de tamaño y sesgo en los datos.

\citet{chatterjee2022multimodal} siguen otra trayectoria: en vez de depender únicamente del texto bruto, construyen una técnica multiatributo que combina rasgos lingüísticos, temáticos, sentimentales, temporales, de personalidad y de comportamiento en Reddit y Twitter. El valor de este enfoque radica en que la mejora predictiva no proviene solamente de una arquitectura más compleja, sino de una representación del problema más rica. Desde la dimensión de efectividad predictiva, esto permite integrar señales complementarias; desde la dimensión de interpretabilidad clínica, introduce variables con mayor potencial explicativo, pues el tiempo de publicación, los emoticonos, los rasgos afectivos y los perfiles de interacción pueden convertirse en señales más legibles que un vector contextual opaco.

La evaluación de severidad propuesta por \citet{gaur2019severity} profundiza todavía más en la conexión entre señal computacional y significado clínico. En lugar de clasificación binaria, los autores formulan un esquema de cinco niveles inspirado en C-SSRS, diferencian usuarios de apoyo de usuarios con diferentes grados de riesgo y utilizan recursos médicos para crear un léxico de severidad suicida. Su CNN logra el mejor equilibrio entre \textit{graded recall}, error ordinal y medida de riesgo percibido. A diferencia de los enfoques puramente predictivos, esta formulación es especialmente valiosa para la interpretabilidad clínica porque organiza la salida del sistema en niveles accionables y no solo en etiquetas abstractas. Este trabajo aporta al reto porque aproxima la clasificación automática a decisiones de triage más cercanas al lenguaje profesional.

Los trabajos más recientes refuerzan la búsqueda de mayor desempeño con arquitecturas profundas. \citet{mirtaheri2024tcn} proponen un modelo LSTM--TCN con autoatención sobre Reddit y Twitter, argumentando que la combinación mejora memoria de largo alcance, escalabilidad y extracción de rasgos relevantes. \citet{shuvo2024bigru} comparan varios modelos de lenguaje preentrenados acoplados con Bi-GRU y reportan un rendimiento muy alto, además de bajas tasas de falsos negativos, aspecto especialmente crítico en aplicaciones de riesgo. Desde la dimensión de efectividad predictiva, ambos artículos refuerzan la frontera cuantitativa del campo. No obstante, ambos muestran una limitación común del estado del arte: aun cuando la precisión aumenta, la explicación clínica de por qué una publicación fue catalogada como riesgosa sigue siendo secundaria.

\subsection{Detección de depresión y enfoques clínicamente fundamentados en DSM-5}

La literatura sobre depresión en redes sociales ha pasado de la detección general del trastorno a esquemas más cercanos a la práctica clínica. \citet{gaur2018dsm5} representan un antecedente importante al mapear subreddits de salud mental a categorías DSM-5 mediante modelado temático, \textit{embeddings} y bases de conocimiento médico. Su propuesta SEDO reduce de manera sustancial la tasa de falsas alarmas y obtiene una concordancia alta con expertos. Desde la dimensión de interpretabilidad clínica, el trabajo es central porque introduce una mediación semántica entre lenguaje social y categorías diagnósticas. Aunque el objeto de análisis son subreddits y no pacientes individuales, el estudio demuestra que la incorporación de conocimiento clínico estructurado puede mejorar simultáneamente precisión y trazabilidad semántica.

Ese movimiento hacia una mayor granularidad clínica se consolida con \citet{bao2025redsm5}, quienes presentan ReDSM5, un corpus de 1484 publicaciones largas de Reddit anotadas a nivel de oración por una psicóloga con los nueve síntomas depresivos del DSM-5 y racionales clínicos breves. Este recurso es uno de los aportes más importantes del corpus revisado para la dimensión de interpretabilidad clínica: la tarea deja de ser ``depresión sí/no'' y pasa a ser identificación multietiqueta de síntomas concretos, con evidencia textual y explicación. Desde la dimensión de efectividad predictiva, los resultados base muestran que los modelos basados en transformers ajustados superan con claridad a SVM y CNN, lo que sugiere que ambas metas pueden coexistir cuando el diseño del dataset lo permite. Este trabajo aporta al reto porque ofrece una referencia directa de cómo vincular señales lingüísticas con síntomas clínicos observables.

\citet{qasim2025depression} complementan esta línea al centrarse en la severidad depresiva y comparar enfoques basados en n-gramas, \textit{sentence transformers} y modelos basados en transformers ajustados de extremo a extremo. Sus resultados muestran que modelos como DeBERTa capturan mejor las diferencias entre clases leve, moderada y severa que las representaciones estáticas o los modelos con ingeniería manual de atributos. Desde la dimensión de efectividad predictiva, el artículo confirma la ventaja de los modelos basados en transformers finamente ajustados. Aun así, mantiene una interpretabilidad limitada en comparación con los trabajos DSM-5, pues la explicación se apoya principalmente en el poder contextual del modelo y no en racionales clínicos explícitos.

\subsection{Aprendizaje profundo, modelos basados en transformers y modelos de lenguaje de gran escala}

La fase más reciente del estado del arte desplaza el foco desde la simple clasificación hacia modelos capaces de representar contexto amplio, extraer evidencia y apoyar interpretación. \citet{bauer2024llm} utilizan \textit{embeddings} de modelos de lenguaje, reducción de dimensionalidad y apoyo de un modelo generativo para analizar 2.9 millones de publicaciones de treinta subreddits. Su objetivo principal no es predecir riesgo individual, sino comprender dimensiones latentes del lenguaje suicida. Desde la dimensión de interpretabilidad clínica, este artículo es especialmente relevante, ya que conecta patrones lingüísticos con teorías contemporáneas del suicidio como desconexión, carga percibida, desesperanza o trauma.

Por su parte, \citet{singh2024evidence} exploran una tarea todavía más cercana al apoyo experto: extraer y resumir evidencia textual de ideación suicida con LLMs. Su desempeño sobresaliente en la tarea compartida de CLPsych 2024 muestra que los modelos generativos pueden producir salidas útiles no solo para clasificar, sino también para justificar. A diferencia de los enfoques puramente predictivos, esta transición sitúa la interpretabilidad en el nivel de fragmentos de evidencia y resúmenes estructurados, no únicamente en variables agregadas o pesos de atención. Este trabajo aporta al reto porque muestra una vía concreta para convertir el resultado del modelo en evidencia comunicable.

En conjunto, los artículos de esta sublínea sugieren una conclusión provisional. Los modelos basados en transformers y los LLMs elevan claramente la capacidad del campo para modelar contexto, severidad y relaciones semánticas complejas; sin embargo, su utilidad clínica depende de cómo se diseñe el problema. Cuando se usan solo para maximizar exactitud, la interpretabilidad sigue siendo débil. Cuando se combinan con anotación experta, taxonomías clínicas o tareas de extracción de evidencia, se convierten en herramientas más pertinentes para salud mental \citep{bao2025redsm5, bauer2024llm, singh2024evidence}.

Esta evolución también muestra que los modelos de lenguaje no sustituyen por completo a los enfoques previos, sino que amplían el tipo de preguntas que pueden formularse. Mientras los modelos clásicos y profundos se concentran principalmente en clasificación, los LLMs permiten explorar tareas de extracción de evidencia, generación de explicaciones y análisis semántico de mayor granularidad. Para el reto, esta diferencia es importante porque sugiere que la comparación futura entre enfoques debe considerar tanto el desempeño numérico como el tipo de salida producida por cada modelo.

\newcolumntype{P}[1]{>{\raggedright\arraybackslash}p{#1}}

\begin{table*}[!t]
\centering
\caption{Corpus analizado para la revisión del estado del arte, organizado por prioridad de lectura y caracterización metodológica.}
\label{tab:state-of-the-art}
\renewcommand{\arraystretch}{1.08}
\setlength{\tabcolsep}{4pt}
\footnotesize
\begin{adjustbox}{max width=\textwidth}
\begin{tabular}{P{0.05\textwidth} P{0.11\textwidth} P{0.12\textwidth} P{0.13\textwidth} P{0.12\textwidth} P{0.08\textwidth} P{0.19\textwidth} P{0.17\textwidth}}
\toprule
Prioridad & Problema & Fuente de datos & Representación textual & Modelo & Métricas & Contribución principal & Limitaciones \\
\midrule
1 & Detección general de trastornos mentales & Literatura científica y múltiples plataformas & Rasgos léxicos, semánticos y metadatos & Revisión de literatura & F1, P, R, Acc., AUC, ROC, ERDE & Sistematiza tendencias, conjuntos de datos, algoritmos y retos del área & No propone un modelo nuevo y depende de evidencia heterogénea \\
2 & Ideación suicida en redes sociales & 37 estudios sobre Twitter, Reddit, Weibo y microblogs & Rasgos lingüísticos, psicométricos, metadatos e interacción & Revisión de métodos de ML y DL & P, R, F1, Acc. & Resume plataformas, conjuntos de datos y algoritmos para suicidalidad & Sesgo al inglés y comparación indirecta entre estudios \\
3 & Depresión basada en síntomas DSM-5 & 1484 publicaciones largas de Reddit & Oraciones anotadas por síntoma y racionales clínicos & SVM, CNN, BERT y LLM ajustado & Acc.; micro, macro y weighted F1 & Presenta ReDSM5 con anotación experta y explicación a nivel síntoma & Cobertura limitada al inglés y alto costo de anotación profesional \\
4 & Detección temprana de ideación suicida & Reddit & N-grams, BoW, TF-IDF y \textit{embeddings} & CNN, LSTM, LSTM-CNN, SVM, NB, RF, XGBoost & Acc., P, R, F1 & Valida que LSTM-CNN supera baselines clásicos & Conjunto de datos limitado y baja interpretabilidad \\
5 & Ideación suicida multimodal & Reddit y Twitter & Rasgos lingüísticos, tópicos, afecto, temporalidad y personalidad & LR, SVM, RF, XGBoost & Acc., P, R, F1, kappa & Integra seis grupos de rasgos para mejorar detección y comparación & Alta ingeniería de atributos y menor generalización \\
6 & Cambio desde discurso de salud mental hacia ideación suicida & Subreddits de salud mental y r/SuicideWatch & Variables lingüísticas, interacción social y tokens de tratamiento & Emparejamiento por propensión y regresión logística & P, R, F1, ROC, AUC & Modela transición temporal y detecta marcadores conductuales de riesgo & Supone continuidad entre cuentas y depende de una ventana temporal específica \\
7 & Mapeo de contenido de Reddit a categorías DSM-5 & 15 subreddits de salud mental & Tópicos, \textit{embeddings} y léxicos médicos & Clasificador multiclase con SEDO y ontologías & TPR, F1, falsas alarmas, kappa & Vincula contenido social con categorías clínicas y reduce falsas alarmas & Analiza subreddits, no usuarios o publicaciones individuales \\
8 & Severidad de riesgo suicida & 2181 usuarios de Reddit; gold standard de 500 & Contenido normalizado con recursos médicos y léxico de severidad & CNN, SVM, RF, FFNN y baseline basado en reglas & Graded precision, graded recall, F-score, error ordinal & Propone un esquema de cinco niveles inspirado en C-SSRS & Conjunto de datos costoso y evaluación ordinal compleja \\
9 & Detección temprana de ideación suicida & 232074 publicaciones de Reddit & Texto preprocesado y \textit{embeddings} contextuales & BERT, RoBERTa, DistilBERT, DistilRoBERTa, ELECTRA + Bi-GRU & Acc., P, R, F1, falsos negativos & Muestra la ventaja de combinar modelos de lenguaje con Bi-GRU & Prioriza rendimiento y ofrece poca explicación clínica \\
10 & Comprensión lingüística de la suicidalidad & 2.9 millones de publicaciones de 30 subreddits de Reddit & \textit{Embeddings} de oraciones y dimensiones latentes & \textit{Embeddings} de LLM, reducción de dimensionalidad y XAI generativa & Análisis interpretativo & Relaciona patrones lingüísticos con teorías del suicidio y taxonomías dimensionales & No valida contra juicio clínico externo y usa una ventana temporal corta \\
11 & Extracción y resumen de evidencia suicida & Dataset UMD/CLPsych con publicaciones de r/SuicideWatch & Texto del post y meta-información de riesgo, emoción y sentimiento & Mixtral 7B$\times$8 y Tulu-2-DPO-70B con \textit{prompting} & F1 de extracción y consistencia media & Convierte la tarea en recuperación de evidencia interpretable con LLMs & Depende del \textit{prompting} y de una tarea compartida acotada \\
12 & Ideación suicida en publicaciones sociales & Twitter y Reddit & Análisis léxico, semántico y secuencial & LSTM + Bi-TCN con autoatención (AL-BTCN) & Acc., P, R, F1 & Propone una arquitectura escalable con memoria de largo alcance & Menor interpretabilidad y necesidad de más validación real \\
13 & Severidad de depresión & Reddit & N-grams, \textit{sentence transformers} y \textit{embeddings} contextuales & XGB, RF, MLP, RoBERTa, BERT, DeBERTa & P, R, F1, Acc. & Compara líneas clásica, contextual y basada en transformers para severidad depresiva & Falta anclaje clínico explícito a síntomas DSM-5 o racionales expertos \\
\bottomrule
\end{tabular}
\end{adjustbox}
\vspace{2pt}

\scriptsize\textit{Nota.} P = precisión, R = \textit{recall}, Acc. = \textit{accuracy}.
\end{table*}

\FloatBarrier
\section{Análisis comparativo}

La Tabla~\ref{tab:state-of-the-art} permite observar una regularidad importante: la literatura revisada no avanza en una sola dirección, sino en dos ejes parcialmente tensionados. El primero es el de \textbf{efectividad predictiva}. En este eje, los trabajos más fuertes son aquellos que explotan arquitecturas profundas o modelos basados en transformers, como LSTM-CNN, Bi-GRU con modelos preentrenados, TCN con autoatención y modelos ajustados para severidad depresiva \citep{tadesse2020suicide, mirtaheri2024tcn, shuvo2024bigru, qasim2025depression}. Estos enfoques suelen reportar mejores resultados cuantitativos porque modelan dependencias largas, contexto y variación semántica con mayor riqueza que las aproximaciones basadas en n-gramas o léxicos aislados.

El segundo eje es el de \textbf{interpretabilidad clínica}. Desde esta dimensión, destacan menos los clasificadores de alto rendimiento y más los trabajos que reestructuran la tarea para alinearla con marcos clínicos o evidencias observables. \citet{gaur2018dsm5} y \citet{gaur2019severity} incorporan taxonomías médicas y esquemas ordinales útiles para intervención; \citet{bao2025redsm5} introducen anotación a nivel de síntoma y racionales; \citet{singh2024evidence} convierten el problema en extracción y resumen de evidencia; y \citet{bauer2024llm} utilizan LLMs para interpretar dimensiones discursivas, no solo para predecir. En otras palabras, la interpretabilidad mejora cuando el diseño metodológico obliga al sistema a producir salidas más cercanas al lenguaje clínico.

También se observan diferencias en la unidad de análisis. Algunos artículos clasifican publicaciones individuales. Otros modelan usuarios, trayectorias o cambios temporales. Un tercer grupo analiza comunidades completas o subreddits como indicadores diagnósticos indirectos \citep{tadesse2020suicide, chatterjee2022multimodal, mirtaheri2024tcn, dechoudhury2016shifts, gaur2019severity, gaur2018dsm5, bauer2024llm}. Esta diferencia no es menor: los enfoques por publicación tienden a maximizar detección inmediata, mientras que los enfoques por usuario o transición aportan mayor valor para monitoreo temprano y contextualización clínica. Este contraste ayuda a explicar por qué el estado del arte no debe leerse como una secuencia de resúmenes, sino como una comparación de decisiones metodológicas con consecuencias distintas.

Por último, la comparación muestra que la calidad de los datos condiciona fuertemente el tipo de conclusión posible. Los artículos de revisión señalan que gran parte del campo depende de etiquetas heurísticas, diccionarios o autodeclaraciones \citep{montejo2024survey, abdulsalam2024suicidal}. En cambio, los trabajos con mayor potencial clínico invierten más en anotación experta y en estructuras diagnósticas explícitas, pero a costa de datasets más pequeños y costosos de construir \citep{gaur2019severity, bao2025redsm5}. La lección metodológica es clara: no existe una sola frontera del estado del arte; la frontera predictiva y la frontera interpretativa avanzan a ritmos distintos y exigen decisiones de diseño diferentes. Este punto resulta especialmente importante para el reto, porque anticipa que una solución sólida deberá justificar tanto su rendimiento como el tipo de evidencia que ofrece.

\section{Brechas y oportunidades de investigación}

La primera brecha es la \textbf{escasez de datos clínicamente anotados}. Los mejores resultados cuantitativos suelen provenir de grandes corpus con etiquetas simplificadas, mientras que los recursos más útiles para interpretación clínica, como los alineados con C-SSRS o DSM-5, son considerablemente más pequeños y costosos de construir \citep{gaur2019severity, bao2025redsm5}. Esto limita la posibilidad de entrenar modelos potentes sin sacrificar granularidad clínica.

La segunda brecha es la \textbf{desconexión entre desempeño y explicabilidad}. Muchos artículos reportan exactitud, F1 o AUC competitivos, pero ofrecen explicaciones débiles sobre qué evidencia sustenta la decisión del modelo. Desde la dimensión de efectividad predictiva, esto puede parecer suficiente; sin embargo, desde la dimensión de interpretabilidad clínica, la limitación es central, porque una predicción útil debería poder relacionarse con síntomas, racionales textuales o indicadores observables para disminuir el riesgo de uso acrítico del sistema \citep{montejo2024survey, mirtaheri2024tcn, shuvo2024bigru}.

La tercera brecha es la \textbf{limitada validación ecológica}. Varios estudios trabajan con Reddit o Twitter como fuente principal y, aunque ello es metodológicamente valioso, todavía existe poca evidencia de generalización hacia otros idiomas, otras plataformas o contextos clínicos reales. Los artículos de revisión coinciden en señalar un sesgo fuerte hacia el inglés y hacia configuraciones experimentales controladas \citep{abdulsalam2024suicidal, montejo2024survey}.

La cuarta brecha es la \textbf{insuficiente integración temporal y multimodal}. El cambio de estado, la progresión de riesgo y la historia previa del usuario son clínicamente relevantes, pero no siempre quedan modelados cuando la tarea se formula como clasificación estática de una sola publicación. Del mismo modo, el campo ha explorado más texto que combinaciones con metadatos, imágenes o patrones de interacción. Varias revisiones identifican precisamente ese camino como una dirección prioritaria \citep{dechoudhury2016shifts, chatterjee2022multimodal, montejo2024survey}.

La quinta brecha es la \textbf{falta de criterios homogéneos para medir interpretabilidad clínica}. Algunos trabajos la asocian con categorías DSM-5, otros con niveles de severidad, otros con recuperación de evidencia y otros con análisis \textit{post hoc} mediante LLMs. Esta variedad es valiosa, pero también dificulta comparar enfoques y seleccionar un marco consistente para futuras fases del proyecto. En consecuencia, una oportunidad clara consiste en diseñar sistemas que combinen potencia predictiva con salidas explícitas y verificables, por ejemplo, síntomas, evidencia textual y nivel de riesgo. Este punto aporta una guía concreta para el reto, porque sugiere que la comparación entre modelos debe considerar no solo quién clasifica mejor, sino también quién explica mejor.

\section{Implicaciones metodológicas para fases posteriores}

A partir del estado del arte revisado, una primera implicación para las siguientes fases del proyecto es la conveniencia de comenzar con un \textit{baseline} reproducible y metodológicamente transparente. Una combinación de TF-IDF con un clasificador lineal constituye una referencia adecuada porque permite establecer un punto de comparación claro, barato de entrenar y fácil de auditar antes de escalar hacia modelos más complejos. Este punto de partida debería integrarse dentro de un \textit{pipeline} modular, de modo que etapas como preprocesamiento, representación textual, clasificación y evaluación puedan sustituirse sin rehacer toda la herramienta.

Una segunda implicación es que la validación futura debe priorizar métricas acordes con el riesgo del problema. ROC AUC conviene como métrica central para comparar modelos, pero debe complementarse con \textit{recall}, F1, precisión y un análisis explícito de falsos negativos, dado que omitir casos de riesgo puede ser más costoso que incurrir en algunas falsas alarmas. Asimismo, si el conjunto de datos disponible incluye \texttt{user\_id}, será necesario evitar fuga de información entre entrenamiento y prueba a nivel de usuario, de modo que el desempeño reportado refleje capacidad real de generalización y no reutilización indirecta de patrones del mismo individuo.

Una tercera implicación es que el diseño conceptual de la herramienta deberá incorporar desde el inicio algún mecanismo básico de interpretabilidad. Esto puede materializarse mediante términos relevantes, n-gramas destacados, pesos del clasificador o fragmentos de evidencia que acompañen la predicción. Estas decisiones metodológicas no se formulan aquí como una solución final propuesta, sino como un conjunto de criterios para orientar la Fase 2A; sobre esa base, será posible comparar posteriormente el \textit{baseline} reproducible con variantes más complejas basadas en \textit{transformers} o LLMs.

\section{Conclusión}

La literatura revisada confirma una evolución clara en la detección computacional de señales de salud mental en redes sociales: el campo ha pasado de enfoques basados en rasgos léxicos y clasificadores convencionales a sistemas profundos, modelos basados en transformers y LLMs con una capacidad predictiva superior. Sin embargo, este avance no resuelve por sí mismo la exigencia central del dominio: producir salidas útiles, justificables y clínicamente comprensibles. Desde la dimensión de efectividad predictiva, el campo ha progresado de forma acelerada. Desde la dimensión de interpretabilidad clínica, en cambio, el progreso ha sido más desigual. Por ello, los trabajos más valiosos para este proyecto no son únicamente los que reportan la mejor exactitud, sino aquellos que articulan el rendimiento con estructuras interpretables, como niveles de severidad, categorías DSM-5, racionales expertos o extracción explícita de evidencia \citep{gaur2018dsm5, gaur2019severity, bao2025redsm5, singh2024evidence}.

En síntesis, el estado del arte sugiere que la dirección más prometedora no consiste en escoger entre desempeño e interpretabilidad, sino en diseñar tareas, datos y modelos que hagan compatibles ambas metas. Reddit se mantiene como una fuente especialmente valiosa por su riqueza discursiva, pero los retos de sesgo lingüístico, generalización, anotación experta y validación clínica siguen abiertos. Para las fases posteriores, esto implica buscar un equilibrio explícito entre desempeño cuantitativo, evidencia textual, modularidad y diseño reproducible, de modo que la herramienta resultante no solo clasifique mejor, sino que también permita justificar y auditar mejor sus decisiones.


\bibliographystyle{elsarticle-harv}
\bibliography{references}

\end{document}
